% ===================================================================
%  toggle_helper.tex - how to write one paper for several versions
%  Version toggles and workflow: Omid Mirzaeedodangeh.
% ===================================================================
%  For reading only: it is never compiled or \input. Lines without %
%  are LaTeX to copy into sections/*.tex; the % lines explain them.
%
%  Versions of this template (change them in setup.tex):
%    conference   versions/conference.tex   IEEE, two columns, 6 pages
%    journal      versions/journal.tex      Elsevier, one column
%    arxiv        versions/arxiv.tex        article, one column, proofs
%    long         a group: journal and arxiv together
%  The review copy (versions/review.tex, or main.tex on Overleaf)
%  shows every version at once, each part in its version's colour.
% ===================================================================


% ===================================================================
%  COMMANDS
% ===================================================================
%  arXiv only                  \onlyin{arxiv}{...}
%                              \begin{arxiv} ... \end{arxiv}
%  conference only             \onlyin{conference}{...}
%                              \begin{conference} ... \end{conference}
%  journal and arXiv           \onlyin{long}{...}
%  not in conference           \exceptin{conference}{...}
%  two versions, by name       \onlyin{journal,arxiv}{...}
%  two wordings                \ifversion{conference}{short}{long}
%  a figure for a group        \ifversion{long}{...}{}
%  a pointer to the appendix   \arxivref{app:proofs}
%  a setting for one version   \ifversion{conference}{...}{}  in setup.tex
%
%  The inline commands take any text, figures and equations included,
%  but no verbatim (\verb, verbatim). The environments, one per version
%  name, take anything; put \begin and \end each on a line of their own.
%  Groups have no environment: use \ifversion{long}{...}{}.


% ===================================================================
%  1. A sentence only in the long versions
% ===================================================================
We compare three controllers%
\onlyin{long}{ and repeat the test with a load step}.
%  conference:  We compare three controllers.
%  journal:     We compare three controllers and repeat the test
%               with a load step.
%  arxiv:       same as journal
%  review copy: the long part in dark blue (shared by two versions)


% ===================================================================
%  2. Two wordings of one sentence
% ===================================================================
We use a \ifversion{conference}{PI controller}{PI controller with anti-windup}.
%  conference:  We use a PI controller.
%  journal:     We use a PI controller with anti-windup.
%  arxiv:       same as journal
%  review copy: both, side by side, with a grey bar between them


% ===================================================================
%  3. The proof only in arXiv, and a pointer to it everywhere
% ===================================================================
%  In the text:
\begin{theorem}\label{thm:stable}
  The closed loop is stable (proof in \arxivref{app:proofs}).
\end{theorem}
%  At the end of main.tex:
\begin{arxiv}
\appendix
\section{Proofs}\label{app:proofs}
% sections/appendix_proofs.tex - the Proofs appendix. main.tex inputs it inside \begin{arxiv}, so only the arXiv version has it.
% For another appendix, add sections/appendix_<name>.tex and input it in main.tex inside \begin{arxiv} ... \end{arxiv},
% To give it to the journal version too, use \ifversion{long}{...}{}; the journal then also needs an \appendix line.
Proofs and extra material go here.

\end{arxiv}
%  conference:  (proof in Appendix A of [5])   and no appendix
%  journal:     (proof in Appendix A of [5])   and no appendix
%  arxiv:       (proof in Appendix A)          and the appendix
%  [5] is the entry with the key arxiv in refs/references.bib: put
%  your arXiv paper there. make -f .github/Makefile copies the labels
%  of the arXiv version to refs/extended_labels.tex. Commit that file
%  so the appendix numbers are right on Overleaf too.


% ===================================================================
%  4. A figure only in the journal and arXiv
% ===================================================================
\ifversion{long}{%
\begin{figure}[t]
  \autofig{example}
  \caption{The speed after a load step.}\label{fig:load}
\end{figure}
}{}
\onlyin{long}{Fig.~\ref{fig:load} shows the load step.}
%  conference:  no figure, no sentence (so no "Fig. ??")
%  journal:     the figure, drawn for one column, and the sentence
%  arxiv:       same as journal


% ===================================================================
%  5. One more citation in the long versions
% ===================================================================
Speed control is often the first loop a student designs~\cite{franklin2019}%
\onlyin{long}{ and~\cite{ogata2010}}.
%  conference:  ... designs [1].
%  journal:     ... designs [1] and [2].
%  Both keys must be in refs/references.bib. Each reference list
%  holds only what that version cites.


% ===================================================================
%  6. Making the conference version fit
% ===================================================================
%  Put each change inside a conference-only command, so the other
%  versions do not change.
%  One line fewer (put it inside the paragraph):
\onlyin{conference}{\looseness=-1}
%  A shorter caption:
\caption{Speed\onlyin{long}{ after a step of the reference}.}
%  Less space under a figure:
\onlyin{conference}{\vspace{-2mm}}
%  A narrower table:
\ifversion{conference}{\setlength{\tabcolsep}{3pt}}{}
%  Tighter floats (in setup.tex):
\ifversion{conference}{\setlength{\textfloatsep}{6pt}}{}
%  The line "% page limit: 6" in versions/conference.tex sets the limit;
%  the laptop build prints OVER when the PDF is longer.


% ===================================================================
%  7. A new version, for example a thesis chapter
% ===================================================================
%  a) In setup.tex add:
\DeclareVersion[Plum]{thesis}
%     and, if it is a long version, change the group line to:
\DeclareVersionGroup{long}{journal,arxiv,thesis}
%  b) Copy versions/arxiv.tex to versions/thesis.tex, change the class
%     line and the page limit, and set \def\paperversion{thesis}.
%  c) Write \onlyin{thesis}{...} or \begin{thesis} ... \end{thesis}.
%  On Overleaf, open versions/thesis.tex and press Recompile.


% ===================================================================
%  8. Co-author edits
% ===================================================================
%  One line per person in setup.tex (replace author1-3):
\DeclareEditor{anna}{AK}   \editorname{anna}{Anna Keller}
\DeclareEditor[RoyalBlue]{ben}{BO}   \editorname{ben}{Ben Osei}
%  In the text:
\added{anna}{new text}
%    review copy: new text in Anna's colour, then AK
%    other PDFs:  new text
\deleted{anna}{old text}
%    review copy: old text struck out     other PDFs: nothing
\replaced{anna}{new}{old}
%    review copy: new, then old struck out     other PDFs: new
\authnote{anna}{add a reference}
%    review copy: [AK: add a reference]     other PDFs: nothing
\begin{addedblock}{anna}
A whole new paragraph.
\end{addedblock}
%    review copy: the paragraph in Anna's colour     other PDFs: the paragraph
\TODO{repeat with ten seeds}
%    review copy: a red box     other PDFs: nothing, and a warning in the log
\pending{0.4~ms}
%    a number still to confirm: pink in every PDF; the submitted PDFs also
%    log a warning
%  Accepting edits:
\acceptedits{anna}
%    from here on Anna's edits are accepted, the others still show;
%    \showedits{anna} shows them again; \acceptalledits accepts all.
%  The files in versions/ (all but review.tex) accept every edit,
%  so the PDFs you submit have no marks.


% ===================================================================
%  9. Figures
% ===================================================================
\autofig{example}
%    figures/src/example.py, drawn at the line width
\autofig[aspect=0.5]{example}
%    height = 0.5 x width
\autofig[height=4cm]{example}
%    a fixed height
\autofig[font=11pt]{example}
%    base size 11 pt: labels 13 pt, ticks and legend 12 pt; do not set it
%    below the body size
\autofig[compact]{example}
%    all figure text at the body size
%  Each version writes the line width at every \autofig to its .afg
%  file, so one \autofig fits every template. After changing a figure's
%  code or its width, run make -f .github/Makefile and commit
%  figures/auto (Overleaf cannot run Python).


% ===================================================================
%  10. Common mistakes
% ===================================================================
%  Verbatim cannot go inside an inline command:
%    WRONG   \onlyin{arxiv}{\verb|x|}  -> error "File ended while scanning"
%    RIGHT   the arxiv environment, \begin and \end on lines of their own
%
%  Words left behind in the other versions:
%    WRONG   see\onlyin{arxiv}{ Fig.~\ref{fig:extra}} for details.
%            -> conference: "see for details."
%    RIGHT   \onlyin{arxiv}{ (see Fig.~\ref{fig:extra})}
%
%  A reference to a label only one version has gives "??":
%    WRONG   see Appendix~\ref{app:proofs}  -> conference: "see Appendix ??"
%    RIGHT   \arxivref{...} for the arXiv appendix, or keep the
%            reference inside the same \onlyin as the label.
%
%  Version names are case sensitive:
%    WRONG   \onlyin{arXiv}{...}  -> error "Unknown version `arXiv'"
%    RIGHT   conference, journal, arxiv, long (as in setup.tex)
%
%  \deleted and \replaced work inside one paragraph. Put braces
%  around \cite, \ref and math inside them:
%    \replaced{anna}{{\cite{a}}}{{\cite{b}}}
%
%  Do not put the same \autofig on both sides of \ifversion: the review
%  copy prints both. Use one \autofig and set its size per version in
%  setup.tex (see the \stepaspect example there).
%
%  Renaming the version arxiv: main.tex and .github/tex/layout.tex
%  use the name too; rename it there as well.
